%-------------------------------------------------------------------------------

% This file is part of code_saturne, a general-purpose CFD tool.
%
% Copyright (C) 1998-2026 EDF S.A.
%
% This program is free software; you can redistribute it and/or modify it under
% the terms of the GNU General Public License as published by the Free Software
% Foundation; either version 2 of the License, or (at your option) any later
% version.
%
% This program is distributed in the hope that it will be useful, but WITHOUT
% ANY WARRANTY; without even the implied warranty of MERCHANTABILITY or FITNESS
% FOR A PARTICULAR PURPOSE.  See the GNU General Public License for more
% details.
%
% You should have received a copy of the GNU General Public License along with
% this program; if not, write to the Free Software Foundation, Inc., 51 Franklin
% Street, Fifth Floor, Boston, MA 02110-1301, USA.

%-------------------------------------------------------------------------------

\subsubsection{Correction step}\label{sec:timstp:correction_step}\label{ap:resopv}

\hypertarget{resopv}{}

The velocity projection (or pressure correction) step enforces mass conservation
on the velocity field computed during the prediction step. The predicted
velocity field $\favre{\vect{u}}$ obtained from the momentum balance
(\S~\ref{ap:cs_velocity_prediction}) relies on an explicit treatment of the
pressure gradient and does not satisfy the continuity equation.

In this step, the pressure increment $\delta p = p^{n+1} - p^n$ is determined by
solving a Poisson-type equation, and is subsequently used to correct both the cell
velocities and the face mass fluxes. In the standard segregated \textit{SIMPLEC}
solver, velocity and pressure are coupled solely through the
unsteady mass operator. (For calculations at large Courant numbers, a pseudo-coupled
reinforced velocity-pressure coupling is also available, as detailed in
\S~\ref{sec:timstp:reinforced_coupling}.)

From a continuous viewpoint, the projection step consists of solving the Poisson
equation:
\begin{equation}
  \dive \left( \Delta t \, \grad \delta p \right) = \dive \left( \rho \favre{\vect{u}} \right) - \Gamma
  \label{eq:pressure_correction:cont_poisson}
\end{equation}
and then updating the velocity field:
\begin{equation}
  \vect{u}^{n+1} = \favre{\vect{u}} - \frac{\Delta t}{\rho} \, \grad \delta p
  \label{eq:pressure_correction:cont_vel_update}
\end{equation}
where $\Gamma$ is an eventual mass source term.

\paragraph{Operator Formulation}
To analyze the fractional step method rigorously, an operator notation is
adopted. The discrete equation of motion at the intermediate prediction level
is written, for each spatial coordinate $\alpha \in \{1, 2, 3\}$, in the form:
\begin{equation}
  \tens{M}_\alpha^n \tens{R}^n \left( \favre{\vect{V}}_\alpha - \vect{V}_\alpha^n \right)
  + \tens{A}_\alpha^n \favre{\vect{V}}_\alpha
  = - \tens{G}_\alpha \vect{P}^n + \vect{S}_\alpha^n + \tens{I}_{s,\alpha} \favre{\vect{V}}_\alpha
  \label{eq:pressure_correction:discrete_qdm}
\end{equation}
where:
\begin{itemize}
\item $\tens{M}_\alpha^n$ is a diagonal matrix of size $\ncell \times \ncell$
  containing the ratio of the cell volume to the local time step:
  $\tens{M}_\alpha^n(i, i) = \vol{\celli} / \Delta t_{\alpha,\celli}^n$,
\item $\tens{R}^n$ is the diagonal density matrix: $\tens{R}^n(i, i) = \rho_\celli^n$,
\item $\favre{\vect{V}}_\alpha$ is the vector of dimension $\ncell$ containing the
  $\alpha$-th component of the predicted velocity field $\favre{\vect{u}}$,
\item $\vect{V}_\alpha^n$ is the vector of dimension $\ncell$ containing the
  $\alpha$-th velocity component at time level $n$,
\item $\tens{A}_\alpha^n$ denotes the convection-diffusion operator (which may be
  non-linear due to slope limiters),
\item $\tens{G}_\alpha$ is the cell-centered gradient operator in direction $\alpha$,
\item $\vect{P}^n$ is the vector of dimension $\ncell$ containing the cell pressure
  at time level $n$,
\item $\vect{S}_\alpha^n$ is the vector of explicit source terms,
\item $\tens{I}_{s,\alpha}$ is the diagonal matrix representing implicit source
  terms acting on velocity component $\alpha$.
\end{itemize}

Rearranging equation~\eqref{eq:pressure_correction:discrete_qdm} yields:
\begin{equation}
  \left( \tens{M}_\alpha^n + \tens{A}_\alpha^n (\tens{R}^n)^{-1}
  - \tens{I}_{s,\alpha} (\tens{R}^n)^{-1} \right) \tens{R}^n \favre{\vect{V}}_\alpha
  = - \tens{G}_\alpha \vect{P}^n + \vect{S}_\alpha^n
  + \tens{M}_\alpha^n \tens{R}^n \vect{V}_\alpha^n
\end{equation}
Defining the linearized momentum operator $\tens{B}_\alpha$:
\begin{equation}
  \tens{B}_\alpha = \tens{M}_\alpha^n + \tens{A}_\alpha^n (\tens{R}^n)^{-1}
  - \tens{I}_{s,\alpha} (\tens{R}^n)^{-1}
\end{equation}
and grouping over the three coordinate directions into block operators:
\begin{equation}
  \tens{B} = \begin{pmatrix} \tens{B}_1 & 0 & 0 \\ 0 & \tens{B}_2 & 0 \\ 0 & 0 & \tens{B}_3 \end{pmatrix},
  \quad
  \tens{G} = \begin{pmatrix} \tens{G}_1 \\ \tens{G}_2 \\ \tens{G}_3 \end{pmatrix},
  \quad
  \favre{\vect{W}} = \tens{R}^n \favre{\vect{V}} = \begin{pmatrix} \rho^n \favre{\vect{V}}_1 \\ \rho^n \favre{\vect{V}}_2 \\ \rho^n \favre{\vect{V}}_3 \end{pmatrix}
\end{equation}
equation~\eqref{eq:pressure_correction:discrete_qdm} is written in compact form:
\begin{equation}
  \tens{B} \favre{\vect{W}} = - \tens{G} \vect{P}^n + \vect{S}^n
  \label{eq:pressure_correction:compact_prediction}
\end{equation}
where $\vect{S}^n$ encompasses explicit sources and the previous time step momentum:
$\vect{S}_\alpha^n + \tens{M}_\alpha^n \tens{R}^n \vect{V}_\alpha^n$.

The momentum equation evaluated at the end of the time step $n+1$ reads:
\begin{equation}
  \tens{B} \vect{W}^{n+1} = - \tens{G} \vect{P}^{n+1} + \vect{S}^n
  \label{eq:pressure_correction:compact_np1}
\end{equation}
Subtracting equation~\eqref{eq:pressure_correction:compact_prediction} from
equation~\eqref{eq:pressure_correction:compact_np1} gives the velocity correction
relation:
\begin{equation}
  \tens{B} \left( \vect{W}^{n+1} - \favre{\vect{W}} \right) = - \tens{G} \left( \vect{P}^{n+1} - \vect{P}^n \right) = - \tens{G} \delta \vect{P}
  \label{eq:pressure_correction:momentum_difference}
\end{equation}
where $\delta \vect{P} = \vect{P}^{n+1} - \vect{P}^n$.

The velocity field at time $n+1$ must satisfy the discrete continuity equation:
\begin{equation}
  \tens{D} \vect{W}^{n+1} = \vect{\Gamma}
  \label{eq:pressure_correction:discrete_continuity}
\end{equation}
where $\tens{D}$ is the discrete divergence operator and $\vect{\Gamma}_i = \int_{\vol{\celli}} \Gamma \, \dd \Omega$.
Applying $\tens{D}$ to equation~\eqref{eq:pressure_correction:momentum_difference} and
substituting equation~\eqref{eq:pressure_correction:discrete_continuity} yields the
discrete Poisson equation for the pressure increment:
\begin{equation}
  \tens{D} \tens{B}^{-1} \tens{G} \, \delta \vect{P} = \tens{D} \favre{\vect{W}} - \vect{\Gamma}
  \label{eq:pressure_correction:operator_poisson}
\end{equation}

\paragraph{Approximation of $\tens{B}^{-1}$ in the Segregated Solver}
Inverting $\tens{B}$ directly is computationally prohibitive because the
convection-diffusion operator couples all mesh cells. In the segregated SIMPLEC
framework, convection, diffusion, and implicit source terms
are neglected in the velocity correction operator, retaining solely the diagonal
unsteady mass contribution:
\begin{equation}
  \tens{B}_\alpha^{-1} \approx \tens{M}_\alpha^{-1} = \text{diag} \left( \frac{\Delta t_{\alpha,\celli}^n}{\vol{\celli}} \right)
\end{equation}
This yields an isotropic equivalent time step at each cell $\celli$:
$\tens{T}_\celli^n = \Delta t_\celli^n \tens{Id}$.

\paragraph{Collocated Finite Volume Discretization and Odd-Even Decoupling}
On collocated meshes, evaluating the operator product $\tens{D} \tens{B}^{-1} \tens{G}$
directly leads to a $2h$-wide stencil on regular Cartesian grids, which decouples
odd and even cells and generates spurious pressure checkerboarding.

To guarantee compact coupling, \CS replaces $\tens{D} \tens{B}^{-1} \tens{G}$ by a
face-based finite volume Laplacian operator $\tens{L}$:
\begin{equation}
  (\tens{L} \, \delta \vect{P})_\celli
  = \sum_{\fij \in \Facei{\celli}} \Delta t_{ij}^n (\grad \delta p)_{\fij} \cdot \vect{S}_{ij}
  + \sum_{\fib \in \Faceb{\celli}} \Delta t_{b}^n (\grad \delta p)_{\fib} \cdot \vect{S}_{b}
  \label{eq:pressure_correction:laplacian_operator}
\end{equation}
where $(\grad \delta p)_{\fij}$ denotes the facet gradient of the pressure increment,
and $\Delta t_{ij}^n$ is the face-interpolated time step.

The facet gradient flux across an interior face $\fij$ is decomposed into a normal direct
contribution and a non-orthogonal reconstruction:
\begin{equation}
  \Delta t_{ij}^n (\grad \delta p)_{\fij} \cdot \vect{S}_{ij}
  = \Delta t_{ij}^n \frac{\delta p_\cellj - \delta p_\celli}{\overline{I'J'}} S_{ij}
  + \frac{\vect{JJ'} - \vect{II'}}{\overline{I'J'}} \cdot \frac{1}{2}
  \left( \Delta t_\celli^n \grad(\delta p)_\celli + \Delta t_\cellj^n \grad(\delta p)_\cellj \right) S_{ij}
  \label{eq:pressure_correction:face_flux_decomp}
\end{equation}
where $I'$ and $J'$ denote the orthogonal projections of cell centers onto the face
normal line passing through the face center $F$.

\paragraph{Rhie \& Chow Filter on Collocated Meshes}
The predicted mass divergence term $\tens{D} \favre{\vect{W}}$ inherits cell-centered
pressure gradients from the prediction step, which can also trigger checkerboard
instabilities on collocated grids.
To ensure strong pressure-velocity coupling, a Rhie \& Chow filter is applied to the
face mass fluxes:
\begin{equation}
  \widetilde{m}_{ij} = \left[ \rho^n \favre{\vect{u}} + \alpha_{\text{Arak}} \Delta t^n \grad(p^n) \right]_{\fij} \cdot \vect{S}_{ij}
  - \alpha_{\text{Arak}} \, \Delta t_{ij}^n (\grad p^n)_{\fij} \cdot \vect{S}_{ij}
  \label{eq:pressure_correction:rhie_chow_interior}
\end{equation}
and on boundary faces $\fib$:
\begin{equation}
  \widetilde{m}_{b} = \left[ \rho^n \favre{\vect{u}} + \alpha_{\text{Arak}} \Delta t^n \grad(p^n) \right]_{\fib} \cdot \vect{S}_{b}
  - \alpha_{\text{Arak}} \, \Delta t_{b}^n (\grad p^n)_{\fib} \cdot \vect{S}_{b}
  \label{eq:pressure_correction:rhie_chow_boundary}
\end{equation}
where $\alpha_{\text{Arak}} \in [0, 1]$ is the Arakawa weighting coefficient
(default value $\alpha_{\text{Arak}} = 1$).
The term inside square brackets is interpolated from cell-centered values, while
$(\grad p^n)_{\fij} \cdot \vect{S}_{ij}$ is evaluated as a facet normal gradient.

\paragraph{Discrete Poisson System and Non-Orthogonality Resolution}
Combining the discrete Laplacian~\eqref{eq:pressure_correction:laplacian_operator}
and the Rhie \& Chow mass divergence yields the Poisson equation:
\begin{equation}
  \sum_{\fij \in \Facei{\celli}} \Delta t_{ij}^n (\grad \delta p)_{\fij} \cdot \vect{S}_{ij}
  + \sum_{\fib \in \Faceb{\celli}} \Delta t_{b}^n (\grad \delta p)_{\fib} \cdot \vect{S}_{b}
  = \sum_{\fij \in \Facei{\celli}} \widetilde{m}_{ij} + \sum_{\fib \in \Faceb{\celli}} \widetilde{m}_{b}
  - \int_{\vol{\celli}} \Gamma \, \dd \Omega
  \label{eq:pressure_correction:assembled_poisson}
\end{equation}

To account for non-orthogonal mesh cells while preserving a sparse, diagonally
dominant linear system, an incremental fixed-point sequence is resolved:
\begin{equation}
  \delta p^{k+1} = \delta p^k + C_{\text{relax}} \, \delta(\delta p)^{k+1}, \quad \text{with } \delta p^0 = 0
\end{equation}
where $k$ is the non-orthogonality iteration index and $C_{\text{relax}} \in (0, 1]$
is a relaxation factor (default $1$).
At each iteration $k+1$, the sub-increment $\delta(\delta p)^{k+1}$ is computed by
solving:
\begin{equation}
  \sum_{\fij \in \Facei{\celli}} \Delta t_{ij}^n \frac{\delta(\delta p)_\celli^{k+1} - \delta(\delta p)_\cellj^{k+1}}{\overline{I'J'}} S_{ij}
  + \sum_{\fib \in \Faceb{\celli}} \Delta t_{b}^n \frac{\delta(\delta p)_\celli^{k+1} - \delta(\delta p)_{b}^{k+1}}{\overline{I'F}} S_{b}
  = r_\celli^k
  \label{eq:pressure_correction:incremental_linear_sys}
\end{equation}
where the right-hand side residual $r_\celli^k$ incorporates the reconstructed
gradient contributions from iteration $k$:
\begin{equation}
  r_\celli^k = \sum_{\fij \in \Facei{\celli}} m_{ij}^k + \sum_{\fib \in \Faceb{\celli}} m_{b}^k - \int_{\vol{\celli}} \Gamma \, \dd \Omega
\end{equation}
with:
\begin{equation}
  m_{ij}^k = \widetilde{m}_{ij} - \Delta t_{ij}^n (\grad \delta p^k)_{\fij} \cdot \vect{S}_{ij}
\end{equation}
Convergence of the non-orthogonality loop is attained when the $L^2$ norm of the
residual vector $\vect{r}^k$ satisfies:
\begin{equation}
  \| \vect{r}^k \|_2 < 10 \, \varepsilon_p \, \mathcal{R}_{\text{norm}}
\end{equation}
where $\varepsilon_p$ is the user-defined pressure solver tolerance and
$\mathcal{R}_{\text{norm}}$ is the normalized mass residual evaluated from the
predicted velocity divergence:
\begin{equation}
  \mathcal{R}_{\text{norm}} = \left\| \dive(\rho \favre{\vect{u}}) - \Gamma \right\|_2
\end{equation}

\paragraph{Velocity and Mass Flux Updates}
Upon convergence (or at the maximum number of sweeps), the fields are updated:
\begin{enumerate}
\item \textbf{Pressure field:}
  \begin{equation}
    p^{n+1} = p^n + \delta p
  \end{equation}
\item \textbf{Cell-centered velocity field:}
  \begin{equation}
    \vect{u}_\celli^{n+1} = \favre{\vect{u}}_\celli - \frac{\Delta t_\celli^n}{\rho_\celli} \grad(\delta p)_\celli
  \end{equation}
\item \textbf{Face mass fluxes:}
  \begin{equation}
    \dot{m}_{ij}^{n+1} = \widetilde{m}_{ij} - \Delta t_{ij}^n (\grad \delta p)_{\fij} \cdot \vect{S}_{ij}
  \end{equation}
  and similarly on boundary faces.
\end{enumerate}

To strictly guarantee global and local mass conservation $\sum_j \dot{m}_{ij}^{n+1} + \sum_k \dot{m}_{b,ik}^{n+1} = \int_{\vol{\celli}} \Gamma \, \dd \Omega$
at the end of the time step, even when non-orthogonality iterations have not fully
converged, the final flux correction uses the non-reconstructed part of the last
increment $\delta(\delta p)^{k_{\text{end}}}$, matching the divergence of the
solved matrix system exactly.
